\section{Conclusions and outlook}
\label{sec:conclusions}

\pvf{} reconstructs primary vertices by mapping reconstructed tracks directly to a
one-dimensional histogram of vertex density and finding its peaks, at an inference
cost that does not grow with pile-up. The two-network pipeline has been merged into
a single end-to-end model, and the intermediate KDE has been removed from training
by warming up the MLP on the histogram target before enabling the U-Net. At low
pile-up the model matches AMVF and exceeds it in vertex--vertex resolution
($0.42$~mm against $0.76$~mm on Run~2 simulation), and is validated on Run~2 and
Run~3 collision data.

The main result of this note is the extension to HL-LHC PU200. After scaling the
model and stabilizing the optimization schedule, and training on the full
$2.74\times10^{6}$-event sample with HL-LHC-specific target widths derived from an
ITk resolution fit, the algorithm reconstructs clean vertices at a yield slightly
above AMVF ($70.5$ against $67.9$ per event) at a lower fake rate ($14.0$ against
$15.2$ per event, falling to $11.3$ with the post-hoc histogram gate), and a tighter
vertex--vertex resolution of $0.29$~mm. It does so at an
inference cost that is flat with pile-up (Section~\ref{sec:timing}): where AMVF pays a
combinatorial cost that grows with the vertex multiplicity the HL-LHC is designed to
increase, \pvf{} pays the same fixed cost per event. This is the operational reason to
pursue the method at Run~4 even where the reconstruction quality is only comparable.

The chain is completed by the graph-network track-to-vertex association
(Section~\ref{sec:gnn}), which turns the found peaks into per-vertex track
lists at negligible cost ($3.8$~ms per event; the optimized full chain runs
in $\sim$11~ms). Judged by the track-based purity criteria applied
identically to both algorithms, the combined \pvf{}+GNN chain delivers a
clean-vertex yield of $71.6\%$ of truth vertices at PU200 against $57.3\%$
for AMVF, at an eighteen-fold lower fake rate, and identifies the
hard-scatter vertex at least as often ($98.1\%$ against $97.5\%$). The
associator operates within three points of its oracle bound on the given
peaks, so the remaining inefficiency belongs to the finder.

Two findings shape the next steps. Capacity and the optimization schedule do not
seem to be the bottleneck: a fivefold-larger model and a sweep over learning rate,
optimizer and loss all reach the same plateau; the limit lies in the
mean-squared-error loss against fixed-width Gaussian targets. A non-negligible
part of the residual inefficiency is physical: about $15\%$ of adjacent truth
vertices at PU200 are closer than the resolution and cannot be separated at the
current output granularity. The fake-versus-real signal is already present in the
predicted histogram, so a learned post-hoc gate recovers most of the available
separation, and an in-model objectness head would mainly serve to clean up the
histogram during training.

The priorities going forward are loss functions that penalise fake,
merged and missed vertices directly rather than through a uniform mean-squared
error, with an in-model objectness head as the concrete first step; a finer output
representation or a deconvolution step to recover the close pairs that lie above the
physical ceiling; and, with the association step now essentially saturated on the
peaks it is given, joint finder--associator objectives that emit separate peaks
wherever the downstream track assignment could distinguish them.
